\subsection{Counterfactual numerical comparison of decision authority}
\label{app:mechanism_numerical}

To complement the analytical distinction in
Section~\ref{subsec:mechanism_comparison}, we conduct a controlled
counterfactual experiment on the same market scale used in the main
simulation. Specifically, all formulations are evaluated for \(T=120\)
monthly rounds with \(N=100\) patients and \(M=4\) insurers using \(50\)
paired random seeds. For each seed, the patient states, private affordability
constraints, coverage requirements and stochastic risk shocks are held common
across all mechanisms. Hence, differences arise from the decision structure
rather than from different sampled populations.

The centralized welfare comparator is permitted to directly replace patient
assignments, insurer premiums, coverage plans and transfers, subject to the
same actuarial cost technology and a minimum insurer participation margin.
The Bayesian comparator uses the reported private type
\(\theta_{i,t}=(B_{i,t},C^{\mathrm{req}}_{i,t})\), a population prior, and
report-dependent allocation and transfer rules. These two implementations are
used as transparent counterfactual representatives of the broader formulations
in Eqs.~\eqref{eq:classical_welfare} and \eqref{eq:bayesian_design}; they are
not intended as replicas of any particular empirical insurance mechanism.

For FL-IMDR, insurer premiums and coverage plans remain generated by the
insurers' response rules and each patient re-solves its own best response after
receiving a voucher. The regulator therefore does not overwrite either class
of decisions. To quantify this distinction, define the endogenous-decision
override residual
\begin{equation}
\begin{aligned}
R_t
=&\;
\frac{1}{N}\sum_{i=1}^{N}
\mathbf 1\!\left\{
j_{i,t}\neq
\operatorname{BR}_{i,t}^{\mathcal P}
(\mathbf p_t,\mathbf c_t,v_{i,t})
\right\}\\
&+\frac{1}{2NM}\sum_{i=1}^{N}\sum_{j=1}^{M}
\left(
\frac{|p_{i,j,t}-p_{i,j,t}^{\mathrm{BR}}|}
{\max(p_{i,j,t}^{\mathrm{BR}},1)}
+
|c_{i,j,t}-c_{i,j,t}^{\mathrm{BR}}|
\right),
\end{aligned}
\label{eq:override_residual}
\end{equation}
where \((p_{i,j,t}^{\mathrm{BR}},c_{i,j,t}^{\mathrm{BR}})\) denote the
offers generated by the autonomous insurer response model. Thus,
\(R_t=0\) indicates that insurer offers are not directly replaced and patient
assignments remain consistent with their decentralized best responses.

Figure~\ref{fig:mechanism_coverage_comparison} reports the insured-population
trajectory under common market realizations, whereas
Fig.~\ref{fig:mechanism_override_comparison} visualizes the structural
difference in Eq.~\eqref{eq:override_residual}. The shaded regions denote
\(95\%\) confidence intervals over paired seeds. The numerical results should
therefore be interpreted as a comparison of regulatory authority and induced
market behavior rather than as a claim that a restricted regulator dominates
an unconstrained social-welfare optimum.

\begin{figure}[t]
\centering
\includegraphics[width=\columnwidth]{Figures/figure_1_insured_trajectory.pdf}
\caption{Insured-population trajectories under the three counterfactual
mechanisms. Curves denote means over paired market seeds and shaded regions
show \(95\%\) confidence intervals.}
\label{fig:mechanism_coverage_comparison}
\end{figure}

\begin{figure}[t]
\centering
\includegraphics[width=\columnwidth]{Figures/figure_2_override_residual.pdf}
\caption{Endogenous-decision override residual in
Eq.~\eqref{eq:override_residual}. FL-IMDR remains at approximately zero
because the regulator modifies vouchers while patient and insurer decisions
remain endogenous responses.}
\label{fig:mechanism_override_comparison}
\end{figure}

% Paste the generated outputs/appendix_results_table.tex here.
